\section{Διάρθρωση του Τεύχους}
\label{sec:intro_structure}

Το υπόλοιπο της παρούσας διπλωματικής εργασίας διαρθρώνεται ως ακολούθως:
\begin{itemize}
    \item Στο \textbf{Κεφάλαιο~\ref{ch:background}}, παρουσιάζεται το θεωρητικό υπόβαθρο, καλύπτοντας τους γενετικούς αλγορίθμους, τα μοντέλα νήσων, τις κατανεμημένες δομές DHT, τα ευρετήρια μάθησης RMI και τις αρχές των μοντέλων surrogate.
    \item Στο \textbf{Κεφάλαιο~\ref{ch:architecture}}, αναλύεται η προτεινόμενη αρχιτεκτονική συστήματος, εστιάζοντας στον αγωγό 3-Tier $\epsilon$-bypass και στο πρωτόκολλο επικοινωνίας.
    \item Στο \textbf{Κεφάλαιο~\ref{ch:implementation}}, περιγράφονται οι λεπτομέρειες υλοποίησης των επιμέρους υποσυστημάτων σε Rust και Python.
    \item Στο \textbf{Κεφάλαιο~\ref{ch:evaluation}}, εκτίθενται τα πειραματικά αποτελέσματα, η ανάλυση επιτάχυνσης (speedup) και οι μετρήσεις κλιμακωσιμότητας.
    \item Στο \textbf{Κεφάλαιο~\ref{ch:conclusion}}, συνοψίζονται τα βασικά συμπεράσματα και διατυπώνονται προοπτικές μελλοντικής έρευνας.
    \item Στο \textbf{Παράρτημα~\ref{ch:appendix}}, παρατίθεται ο αναλυτικός ψευδοκώδικας των κρίσιμων αλγορίθμων του συστήματος.
\end{itemize}
